% =====================================================================
%  The Sayan Universal Gradient (SUG) Theory
%  arXiv preprint -- physics.gen-ph
%
%  Compile with:  pdflatex sug_theory_arxiv.tex  (run twice)
%  or:            latexmk -pdf sug_theory_arxiv.tex
%
%  Optional revtex4-2 (journal) alternative: replace the \documentclass
%  and the geometry block below by
%      \documentclass[aps,pre,reprint]{revtex4-2}
%      \usepackage{amsmath,amssymb,amsthm,bm,booktabs,graphicx}
%      \usepackage[usenaturelybrackets]{natbib}
%  and drop the \maketitle block (revtex supplies its own).
% =====================================================================
\documentclass[11pt,a4paper]{article}

% ---- encoding / fonts -------------------------------------------------
\usepackage[T1]{fontenc}
\usepackage[utf8]{inputenc}

% ---- page geometry ----------------------------------------------------
\usepackage[a4paper,margin=1in]{geometry}

% ---- mathematics ------------------------------------------------------
\usepackage{amsmath}
\usepackage{amssymb}
\usepackage{amsthm}
\usepackage{bm}

% ---- tables and figures -----------------------------------------------
\usepackage{booktabs}
\usepackage{graphicx}
\graphicspath{{../figures/}{figures/}{./}}

% ---- links ------------------------------------------------------------
\usepackage[colorlinks=true,linkcolor=blue,citecolor=blue,urlcolor=blue]{hyperref}

% ---------------------------------------------------------------------
%  Notation shortcuts
% ---------------------------------------------------------------------
\newcommand{\thetaS}{\theta_{S}}
\newcommand{\ThetaS}{\Theta_{S}}
\newcommand{\dd}{\mathrm{d}}
\newcommand{\tgrad}{\tau}
\newcommand{\logint}{\operatorname{li}}
\newcommand{\Ei}{\operatorname{Ei}}
\newcommand{\Euler}{\gamma_{\mathrm{E}}}
\newcommand{\code}[1]{\texttt{#1}}

% Framed call-outs (core LaTeX only -- no extra packages required)
\newcommand{\statusnote}[1]{%
  \begin{center}%
  \fbox{\begin{minipage}{0.92\textwidth}\small #1\end{minipage}}%
  \end{center}}

\newcommand{\keyresult}[1]{%
  \begin{center}%
  \fbox{\begin{minipage}{0.92\textwidth}#1\end{minipage}}%
  \end{center}}

% \IfFileExists-guarded figure: compilation never fails on a missing file
\newcommand{\sugfigure}[3]{%
  \begin{figure}[tb]
    \centering
    \IfFileExists{#1}%
      {\includegraphics[width=0.92\linewidth]{#1}}%
      {\fbox{\parbox[c][3.5cm][c]{0.92\linewidth}{\centering
          \textit{Figure file \code{#1} not found in the current path.%
          See the Code Availability section.}}}}%
    \caption{#3}
    \label{fig:#2}
  \end{figure}}

% ---------------------------------------------------------------------
\title{\bfseries The Sayan Universal Gradient (SUG) Theory:\\
       A Proposed Framework Coupling Energy Release, Thermal Gradients
       and Distance}
\author{Sayan M.}
\date{Preprint, version 2.1 (2026)}

\begin{document}
\maketitle

\begin{abstract}
We introduce the \emph{Sayan Universal Gradient (SUG)} theory: an original
and deliberately minimal proposal in which a single dimensionless scalar
field, $\thetaS$, organises three observables of a high-energy source --- the
released energy potential $E_{nb}$, the local thermal gradient $\Delta T$, and
the radial distance $r$ from the source. All physical inputs are normalised
against stated reference scales, giving dimensionless ratios $e$, $\tgrad$ and $x$,
and the field then falls off only \emph{logarithmically} with distance,
\begin{equation}
  \thetaS(x) \;=\; K\,\frac{\sqrt{e}}{\tgrad\,\ln x},
  \qquad
  e=\frac{E_{nb}}{E_0},\quad \tgrad=\frac{\Delta T}{T_0},\quad
  x=\frac{r}{r_0},\quad x>1 ,
\end{equation}
where $K$ is the single dimensionless calibration constant of the model. Five
postulates anchor the framework (energy coupling, logarithmic decay,
square-root energy response, thermal antagonism, and a material threshold
separating bound matter from an information/energy degree of freedom). From
these we derive closed-form expressions for the radial path-integrated
(contour) field, for a vacuum destruction radius $R_d$, and for a breaking
energy $E_c$, and we show that the framework is internally consistent under
round-trip inversion. We then state six falsifiable predictions. The central
one (P2) is a \emph{parameter-free} energy ratio,
\begin{equation}
  \frac{E_2}{E_1} = \left(\frac{\ln x_2}{\ln x_1}\right)^{\!2},
\end{equation}
which is independent of the calibration constant $K$, of the material
threshold $\theta_c$ and of the choice of reference scales, and which is
therefore the cleanest experimental handle on the theory. We further show
that this parameter-freeness is a generic property of the framework rather
than a peculiarity of the logarithmic ansatz, and we describe an explicit
least-squares procedure for fixing $K$ from measured events. SUG is offered
as a testable research programme awaiting its first confrontation with data,
not as an established law of nature.
\end{abstract}

% =====================================================================
\section{Introduction}
\label{sec:intro}

Classical treatments of high-energy events keep three ledgers largely apart:
how much energy was released, what thermal gradient is felt at a given place,
and at what distance observable effects persist. This separation is
convenient, but it may be concealing a simpler underlying order. SUG asks a
single deliberately unifying question:

\begin{quote}
\emph{Can one scalar field organise a source's released energy, the local
thermal gradient, and the distance at which effects persist --- and if so,
what geometric laws follow?}
\end{quote}

The framework proposed here answers affirmatively and works the consequences
out in closed form. Its signature --- and its most aggressive, most testable
assumption --- is \emph{logarithmic} decay with distance. A logarithm falls
far more slowly than any power law, which makes the ansatz easy to falsify:
if a measured energy-versus-range curve is not consistent with a purely
logarithmic envelope, the framework is wrong in a way that cannot be
attributed to a nuisance parameter.

\statusnote{%
\textbf{Status of this work.}
SUG is an \emph{original, speculative proposal}. It has not been peer
reviewed, has not been published in a refereed journal, and has not been
validated against any experiment or observation. Until the free calibration
constant $K$ is fitted to a real event, the theory predicts \emph{functional
shapes and ratios}, not absolute numbers. Every limitation is catalogued in
Sec.~\ref{sec:limitations} and in the accompanying open-questions document,
and the code that implements the algebra is public so that the claims can be
checked rather than taken on trust.}

The remainder of this paper is organised as follows. Sec.~\ref{sec:framework}
gives the postulates, the dimensionally honest formulation, the central
operator, its radial gradient, its contour (path-integrated) form, and the
threshold inversions. Sec.~\ref{sec:predictions} states the six falsifiable
predictions, with particular emphasis on the parameter-free ratio P2.
Sec.~\ref{sec:calibration} sets out the empirical calibration of $K$ and the
statistical machinery required to test the framework. Sec.~\ref{sec:conclusion}
concludes, and Sec.~\ref{sec:code} documents the accompanying software.

% =====================================================================
\section{The Sayan Universal Gradient framework}
\label{sec:framework}

\subsection{The five postulates}

SUG is an axiomatic framework. It starts from five premises, stated so that
each can be criticised and, where possible, attacked by data
(Tab.~\ref{tab:postulates}).

\begin{table}[tb]
  \centering
  \caption{The five SUG postulates. None of these is claimed to be derived
  from an established theory; they are the axioms the framework adopts.}
  \label{tab:postulates}
  \begin{tabular}{@{}p{1.5cm}p{5.1cm}p{7.4cm}@{}}
    \toprule
    \textbf{Postulate} & \textbf{Statement} & \textbf{Consequence} \\
    \midrule
    P-I  & A single scalar field $\thetaS$ couples released energy, thermal
           gradient and distance. &
           Thermal, energetic and geometric information about an event is one
           continuous object, not three independent lists. \\[2pt]
    P-II & $\thetaS \propto 1/\ln x$: the field falls with the
           \emph{logarithm} of the dimensionless distance. &
           The predicted reach of a disturbance in vacuum is far longer than
           an inverse-square model suggests. This is the theory's most
           aggressive and most testable assumption. \\[2pt]
    P-III & $\thetaS \propto \sqrt{E_{nb}}$: energy enters through its
           square root. &
           Compressive energy-to-field mapping; a mild saturation, so that
           past a point extra energy buys little extra field. \\[2pt]
    P-IV & $\thetaS \propto 1/\Delta T$: a larger thermal gradient
           \emph{antagonises} the field. &
           Cold points feel a stronger field than warm ones; the theory can
           say \emph{where} along a temperature profile the effect is
           strongest. \\[2pt]
    P-V  & A material threshold $\theta_c$ separates bound matter from an
           information/energy degree of freedom. &
           With a threshold the framework describes \emph{events}, not merely
           a field, and predicts a destruction radius and a breaking energy. \\
    \bottomrule
  \end{tabular}
\end{table}

Postulates P-II, P-III and P-IV fix the functional dependence of the field,
while P-I fixes what the field is a function of. Postulate P-V supplies the
inversion that converts the field into observable geometry. Nothing further
is assumed; every result below is derived from these five statements plus
dimensional consistency.

\subsection{Reference scales and dimensionless ratios}
\label{sec:scales}

A field that couples quantities with different physical dimensions cannot be
dimensionally consistent unless its arguments are themselves dimensionless. We
therefore declare three reference scales --- a reference energy $E_0$, a
reference thermal gradient $T_0$ and a reference length $r_0$ --- which fix
the units of the theory, and we work exclusively with the dimensionless
ratios
\begin{equation}
  e \;=\; \frac{E_{nb}}{E_0}, \qquad
  \tgrad \;=\; \frac{\Delta T}{T_0}, \qquad
  x \;=\; \frac{r}{r_0}.
  \label{eq:ratios}
\end{equation}
This is the standard non-dimensionalisation of the Buckingham $\pi$ family
\cite{buckingham1915}: a problem with three dimensionful inputs and one
dimensionless output admits a unique set of $\pi$-groups up to normalisation,
and \code{(e, tr, x)} is such a choice. The reference scales are \emph{units},
not physics; changing them changes $e$, $\tgrad$, $x$ simultaneously and leaves
every prediction of the theory invariant (verified numerically in
Sec.~\ref{sec:calibration}).

The motivation for this step is not cosmetic. The naive dimensional reading
\begin{equation}
  \ThetaS = \frac{\sqrt{E_{nb}}}{\Delta T\,\log_{10} r}
  \label{eq:naive}
\end{equation}
is ill-formed twice over: $\sqrt{\text{energy}}/\text{temperature}$ does not
cancel, and the logarithm of a dimensionful length is formally undefined.
Declaring reference scales repairs both defects. Equation~\eqref{eq:naive} is
recovered as the special case $K=1$, $E_0=T_0=r_0=1$ with base-$10$
logarithms; it is a convenient laboratory illustration of the framework, not
a separate theory, and we retain it in the software only for historical
reproducibility.

\subsection{The central operator}
\label{sec:operator}

Postulates P-II, P-III and P-IV together with the requirement that
$\thetaS$ be dimensionless select, up to a multiplicative constant, the form
\begin{equation}
  \boxed{\; \thetaS(x) \;=\; K\,\frac{\sqrt{e}}{\tgrad\,\ln x} , \qquad x > 1 \;}
  \label{eq:theta}
\end{equation}
or, in terms of the physical inputs of Eq.~\eqref{eq:ratios},
\begin{equation}
  \thetaS(r) \;=\;
  K\,\frac{\sqrt{E_{nb}/E_0}}
  {\bigl(\Delta T/T_0\bigr)\,\ln\!\bigl(r/r_0\bigr)} .
  \label{eq:theta_phys}
\end{equation}
Every quantity in Eq.~\eqref{eq:theta_phys} is a pure number, so the field is
manifestly dimensionless. It is positive for $K>0$, $\sqrt{e}>0$, $\tgrad>0$ and
$x>1$; we henceforth discuss its magnitude $\lvert\thetaS\rvert$.

Two features deserve comment.

\emph{Inner boundary.} The point $x=1$ (that is, $r=r_0$) is a genuine
singularity, since $\ln 1 = 0$. We interpret it as the edge of the
\emph{source core}: inside $r\le r_0$ the continuum field description is not
defined, and the field description begins where the source ends. This is a
statement about the limits of the model rather than a defect in it, but it
does require a physical prescription for $r_0$ that we do not yet have
(limit O3, Sec.~\ref{sec:limitations}).

\emph{Slow fall-off.} Because $\ln x$ grows without bound, $\thetaS$ decays
more slowly than any inverse power of distance. For example, doubling the
field decay length from $x=10$ to $x=100$ changes the field by a factor of
only $\ln 100/\ln 10 = 2$ at fixed $e$ and $\tgrad$. This is the origin of the
"long reach" claims of the theory, and also of its tension with the
inverse-square law (Sec.~\ref{sec:limitations}).

\subsection{Radial gradient of the field}
\label{sec:gradient}

The local rate of decay of the field along the radial direction is
\begin{equation}
  \frac{\dd\thetaS}{\dd r}
  \;=\; -\,\frac{K\sqrt{e}}{\tgrad\,r_0\,x\,\bigl(\ln x\bigr)^{2}}
  \;=\; -\,\frac{\thetaS(r)}{r\,\ln\!\bigl(r/r_0\bigr)} ,
  \label{eq:grad}
\end{equation}
with $x=r/r_0$. Equation~\eqref{eq:grad} makes the slow fall-off explicit:
the gradient is suppressed by $\ln x$ relative to the field itself, and by a
further factor $r$ from geometric spreading. A useful diagnostic is the local
\emph{logarithmic} decay length
\begin{equation}
  \ell_{\ln}(r) \;\equiv\; -\,\frac{r}{\dd\ln\thetaS/\dd\ln r}
  \;=\; r\,\ln x ,
\end{equation}
which diverges as $r\to\infty$: SUG contains no fixed length scale, and
consequently no asymptotic regime in which the decay becomes
approximate-power-law. This is a sharp qualitative difference from
inverse-square spreading, and it is directly measurable in principle.

\subsection{Contour (path-integrated) form}
\label{sec:contour}

A global version of the field is obtained by integrating the point field along
a radial segment $[r_a, r_b]$ with $r_b > r_a > r_0$. Writing
$r = r_0 x$, so that $\dd r = r_0\,\dd x$,
\begin{equation}
  \ThetaS(a\!\to\!b)
  \;=\; \int_{r_a}^{r_b}\!
  \frac{\sqrt{e}}{\tgrad\,\ln(r/r_0)}\,\dd r
  \;=\; \frac{K\,r_0\,\sqrt{e}}{\tgrad}
  \Bigl[\logint(x_b)-\logint(x_a)\Bigr],
  \label{eq:contour}
\end{equation}
where the logarithmic integral is
\begin{equation}
  \logint(x) \;\equiv\; \int_0^{x}\frac{\dd t}{\ln t}
  \;=\; \Ei(\ln x)
  \;=\; \Euler + \ln\!\ln x
     + \sum_{k=1}^{\infty}\frac{(\ln x)^{k}}{k\,k!},
  \qquad x>1 .
  \label{eq:li}
\end{equation}
Equation~\eqref{eq:li} is the classical logarithmic integral
\cite{abramowitz1964}, implemented exactly via the exponential integral
\code{expi} in the accompanying software.
If path length is instead measured in dimensionless units
($\dd\ell' = \dd x = \dd r/r_0$), the same result appears with the prefactor
$K\sqrt{e}/\tgrad$ alone; the two conventions differ only by the fixed factor
$r_0$, and the implementation reports the former.

Two remarks. First, $\logint$ is monotonically increasing on $(1,\infty)$, so
$\ThetaS(a\!\to\!b) > 0$ and grows with the path length: the contour constant
is a positive, saturating measure of "field accumulated along a ray".
Second, the closed form~\eqref{eq:contour} is exact; the software verifies it
against adaptive quadrature to machine precision, which is a useful check
that the algebra and the code agree.

\subsection{Threshold inversion: geometry and energy}
\label{sec:inversion}

Postulate P-V supplies a material threshold $\theta_c > 0$. Inverting
Eq.~\eqref{eq:theta} at fixed $e$ and $\tgrad$ fixes the geometry; inverting it
at fixed $\tgrad$ and $x$ fixes the energy.

\paragraph{Vacuum destruction radius.}
Setting $\thetaS(x_c)=\theta_c$ gives
\begin{equation}
  \ln x_c = \frac{K\sqrt{e}}{\tgrad\,\theta_c}
  \quad\Longrightarrow\quad
  \boxed{\; x_c = \exp\!\left(\frac{K\sqrt{e}}{\tgrad\,\theta_c}\right),
  \qquad R_d \equiv r_0\,x_c \;}
  \label{eq:radius}
\end{equation}
where $R_d$ is the radius beyond which the source can no longer hold matter
above its integrity threshold. Note the structure of Eq.~\eqref{eq:radius}:
a logarithmically decaying field inverts to an \emph{exponentially} growing
radius in $\sqrt{e}/(\tgrad\theta_c)$. This is the sharpest quantitative
statement the theory makes, and the most exposed to falsification.

\paragraph{Breaking energy.}
Conversely, at fixed $\tgrad$ and $x$,
\begin{equation}
  e_c = \left(\frac{\tgrad\,\theta_c\,\ln x}{K}\right)^{\!2}
  \quad\Longrightarrow\quad
  \boxed{\; E_c \;=\; E_0
    \left(\frac{\tgrad\,\theta_c\,\ln x}{K}\right)^{\!2} \;}
  \label{eq:energy}
\end{equation}
the source energy required to carry a threshold to distance $r$. Since the
required energy grows only quadratically in $\ln x$, extending an effect
outward is predicted to be remarkably cheap; this is the quantitative content
of the range-extension claim tested in Sec.~\ref{sec:predictions}.

\paragraph{Thermal-profile inversion.}
A third inversion is useful in the field: for a prescribed field value
$\theta_\ast$ at radius $r$, the local thermal gradient must be
\begin{equation}
  \Delta T(r) \;=\; T_0\,\frac{K\sqrt{e}}{\theta_\ast\,\ln(r/r_0)} ,
  \label{eq:thermalinv}
\end{equation}
i.e. the thermal gradient \emph{map} along a ray is fixed once the source
energy is given. This is what allows the framework to make a statement about
\emph{where} along a temperature profile a threshold is reached, and it is
implemented by a Brent root-solver for the general case of an
arbitrarily varying $\Delta T(r)$.

\subsection{Consistency identities}
\label{sec:consistency}

Because the inversions above are algebraic rearrangements of
Eq.~\eqref{eq:theta}, two round-trip identities must hold identically --- the
field evaluated at its own critical radius, Eq.~\eqref{eq:rt1}, and the field
produced by the breaking energy at the target distance,
Eq.~\eqref{eq:rt2} --- and the test suite enforces both:
\begin{align}
  \thetaS\!\bigl(x_c\bigr) &= \theta_c ,
  \label{eq:rt1}\\
  \thetaS(e_c,\,x) &= \theta_c .
  \label{eq:rt2}
\end{align}
These are not empirical statements; they are internal-consistency checks
whose purpose is to certify that the software implements the mathematics
faithfully. A failure of either would indicate an implementation error, not a
failure of the theory.

\subsection{Relation to established physics}
\label{sec:context}

SUG is not presented in a vacuum. Table~\ref{tab:anchors} lists the nearest
classical anchor for each ingredient, together with the tension each
creates. The last two rows are the ones a critic should focus on: the
logarithmic ansatz conflicts with the well-established inverse-square spread of
radiative flux, and the square-root energy response is asserted rather than
derived. We prefer to state these plainly than to obscure them.

\begin{table}[tb]
  \centering
  \caption{SUG elements, their classical anchors, and the resulting tension.
  The anchors provide context for comparison; they do not constitute
  validation of SUG.}
  \label{tab:anchors}
  \begin{tabular}{@{}p{3.9cm}p{4.0cm}p{6.7cm}@{}}
    \toprule
    \textbf{SUG element} & \textbf{Classical anchor} & \textbf{Relation / tension} \\
    \midrule
    Energy as a field driver
      & Relativistic energy; radiation flux \cite{einstein1905}
      & SUG uses $\sqrt{E}$ and a logarithmic envelope, not $E/r^{2}$. \\
    Thermal antagonism
      & Thermodynamics, heat transport \cite{clausius1865}
      & The \emph{direction} of the coupling must be tested; standard cooling
        laws do not obviously oppose a field. \\
    Threshold and phase change
      & Critical phenomena, phase transitions
      & Postulate P-V is the phase-transition analogue; $\theta_c$ has no
        microphysical formula. \\
    Logarithmic decay
      & Coulomb and gravitational log potentials in $2$D
      & Must be reconciled with the $3$D inverse-square law, or reinterpreted
        as acting on a derived rather than a raw observable. \\
    Radial scaling
      & Self-similar blast solutions \cite{sedov1959}
      & Classical detonation scaling predicts power-law, not logarithmic,
        range-extension. \\
    \bottomrule
  \end{tabular}
\end{table}

% =====================================================================
\section{Falsifiable predictions}
\label{sec:predictions}

A framework earns its keep only through the statements it dares to make and
the routes by which it could be shown wrong. We list six, each phrased so that
in principle an experiment or observation could confirm or reject it; they are
summarised in Table~\ref{tab:predictions}. Until $K$ is calibrated,
predictions P1, P3--P5 constrain \emph{shapes}; P2 is independent of $K$ and
can be tested immediately.

\subsection{P1 --- Shape of the reach law}

\emph{Statement.} For fixed $\tgrad$ and $\theta_c$, the threshold distance obeys
\begin{equation}
  \ln R_d \;\propto\; \sqrt{E_{nb}} .
  \label{eq:p1}
\end{equation}
Equivalently, $\ln R_d$ plotted against $\sqrt{E_{nb}}$ must be a straight
line. A plot of $\ln R_d$ against $E_{nb}$ or against $\ln E_{nb}$ would
\emph{not} be linear, so the distinction is visible without any fit quality
argument: P1 admits a purely visual falsification.

\subsection{P2 --- The parameter-free energy ratio (central result)}
\label{sec:p2}

This is the cleanest test the theory offers and the one we would run first.

\emph{Statement.} At a fixed thermal gradient, moving a threshold from
distance $x_1$ to $x_2$ costs released energy in the ratio
\begin{equation}
  \boxed{\; \frac{E_2}{E_1}
  \;=\; \left(\frac{\ln x_2}{\ln x_1}\right)^{\!2} \;}
  \label{eq:p2}
\end{equation}
and this ratio is \emph{independent of} $K$, of the material threshold
$\theta_c$, of the reference scales $E_0,T_0,r_0$, and of the value of $\tgrad$
itself (provided $\tgrad$ is common to the two events). To see the cancellations
explicitly, apply Eq.~\eqref{eq:energy} to both events:
\begin{equation}
  \frac{E_2}{E_1}
  = \frac{E_0\!\left(\dfrac{\tgrad\theta_c\ln x_2}{K}\right)^{\!2}}
         {E_0\!\left(\dfrac{\tgrad\theta_c\ln x_1}{K}\right)^{\!2}}
  = \left(\frac{\ln x_2}{\ln x_1}\right)^{\!2},
  \label{eq:p2proof}
\end{equation}
in which $E_0$, $\tgrad$, $\theta_c$ and $K$ all cancel, as
Eq.~\eqref{eq:p2proof} makes explicit. The prediction is thus a statement
about the \emph{functional form} of the decay law alone, which is exactly why
it can be tested before the theory has been calibrated.

\keyresult{%
\textbf{Why P2 is the decisive test.} Equation~\eqref{eq:p2} contains no free
parameter of the model. It can therefore be used (i.\,e.) on a
\emph{second} event, at a different range, without any fitting whatsoever.
If a measured range extension demands a substantially larger energy
multiplier than Eq.~\eqref{eq:p2} predicts, the logarithmic ansatz is
falsified --- no choice of $K$, $\theta_c$ or of units can rescue it.}

\paragraph{Illustration.} The cost of extending the threshold by one decade
$x\mapsto 10x$ is
\begin{equation}
  \frac{E(10x)}{E(x)} = \left(1+\frac{\ln 10}{\ln x}\right)^{\!2},
  \label{eq:decade}
\end{equation}
which falls as $x$ grows. Numerically,
\begin{center}
\begin{tabular}{@{}lrrr@{}}
  \toprule
  Range extension $x_1\to x_2$ & $x_1\to x_2$ & Predicted $E_2/E_1$ \\
  \midrule
  One decade, early        & $10 \to 100$     & $4.00$ \\
  One decade, later        & $100 \to 1000$   & $2.25$ \\
  One decade, far          & $10^{3} \to 10^{4}$ & $1.78$ \\
  Two decades              & $10^{2} \to 10^{6}$ & $9.00$ \\
  \bottomrule
\end{tabular}
\end{center}
The qualitative signature is that the energy cost of a fixed multiplicative
range extension \emph{decreases} as the range grows, approaching
Eq.~\eqref{eq:p2} with $E_2/E_1\to 1$ in the far field. A model whose reach
is set by geometric spreading (an inverse-square law) instead demands a
\emph{constant} factor per decade, $10^{2p}$ for a $1/r^{2p}$ law, which grows
without bound as the range grows. The two behaviours are qualitatively
distinct and therefore easy to separate observationally.

\paragraph{A remark on genericity.} The parameter-freeness of P2 is not an
accident of the logarithmic choice. Write the framework generically as
\begin{equation}
  \thetaS \;=\; \frac{K\,g(e)}{\tgrad\,f(x)},
  \label{eq:generic}
\end{equation}
with $g(e)=\sqrt{e}$ from P-III. Since $\theta_c = K\sqrt{e}/(\tgrad f(x_c))$,
the breaking energy at any distance is $E_c \propto \bigl[f(x)\bigr]^{2}$, and
therefore
\begin{equation}
  \frac{E_2}{E_1} \;=\; \left[\frac{f(x_2)}{f(x_1)}\right]^{\!2},
  \label{eq:genericratio}
\end{equation}
for \emph{any} monotone decay law $f$. Equation~\eqref{eq:genericratio} is the
genuinely parameter-free core of the framework; SUG's specific and
\emph{sharp} claim is that $f=\ln$, so that the right-hand side becomes
Eq.~\eqref{eq:p2}. An experiment that measures $E_2/E_1$ at two or more
separated ranges thereby measures $f$ itself, up to the unobservable
normalisation of the reference length. This also shows how P2 discriminates
against competing models: for an inverse-square law $f \propto x^{2}$ the
measured ratio would be $(x_2/x_1)^{4}$, enormously larger than
Eq.~\eqref{eq:p2} for the same ranges.

\subsection{P3 --- Thermal antagonism}

\emph{Statement.} For fixed source energy and distance,
\begin{equation}
  \thetaS \;\propto\; \frac{1}{\tgrad} .
  \label{eq:p3}
\end{equation}
\emph{Test.} Hold the source fixed and raise the ambient thermal gradient; the
observed disturbance strength at a fixed point should fall hyperbolically. A
measurement that is flat in temperature, or that rises, rejects the coupling
\emph{direction} of P-IV. Note that this prediction, unlike P2, requires
knowledge of $\tgrad$ in absolute units and hence of $T_0$.

\subsection{P4 --- Square-root saturation}

\emph{Statement.} Doubling the released energy multiplies the field by
exactly $\sqrt{2}$, and multiplies the destruction radius by
\begin{equation}
  \frac{R_d(2e)}{R_d(e)} = \exp\!\left[
  \frac{K\sqrt{e}}{\tgrad\,\theta_c}\bigl(\sqrt{2}-1\bigr)\right] ,
  \label{eq:p4}
\end{equation}
which is modest for moderate sources. Field-versus-energy data should
therefore \emph{saturate}: a concave curve, not a straight line. This is a
distinctive, purely qualitative signature and is the easiest of the six to
test in principle, since it requires only the shape of a single curve.

\subsection{P5 --- Vacuum/atmosphere asymmetry}

\emph{Statement.} Because the SUG field decays only logarithmically, an
attenuating medium (atmosphere, circumstellar envelope) that imposes an
additional factor $\exp(-\mu x)$ on top of the field will suppress it far
more severely than vacuum. The framework therefore predicts that a blast in
deep vacuum affects a much larger volume than an identical blast in an
atmosphere --- a statement with no free parameters in its limit and a
single measurable attenuation length $\mu$ otherwise.

\subsection{P6 --- One organising field, not many channels}

\emph{Statement.} Thermal and mechanical signatures of a high-energy event
should be spatially continuous and \emph{jointly} organised by a single
functional form, with the thermal and mechanical channels correlated as
Eq.~\eqref{eq:theta} requires, rather than described by independent radial
profiles per channel. Falsified if the channels decorrelate, or if one of
them is well described while the other is not.

\subsection{Summary and falsification routes}

\begin{table}[tb]
  \centering
  \caption{The six predictions, whether they are free of model parameters,
  and the most direct route to falsifying each.}
  \label{tab:predictions}
  \begin{tabular}{@{}p{0.7cm}p{4.5cm}p{2.4cm}p{6.5cm}@{}}
    \toprule
    \textbf{\#} & \textbf{Prediction} & \textbf{Parameter-free?} &
    \textbf{Most direct falsification} \\
    \midrule
    P1 & $\ln R_d \propto \sqrt{E}$ linear (Eq.~\eqref{eq:p1}) & no &
        plot $\ln R_d$ vs.\ $\sqrt{E}$; curvature kills it \\
    P2 & $E_2/E_1 = (\ln x_2/\ln x_1)^2$ (Eq.~\eqref{eq:p2}) & \textbf{yes} &
        a much larger energy multiplier is required for the same range \\
    P3 & $\thetaS \propto 1/\tgrad$ (Eq.~\eqref{eq:p3}) & no &
        raise ambient temperature; a flat field rejects it \\
    P4 & field saturates with $E$ (Eq.~\eqref{eq:p4}) & no &
        field-vs-energy curve must be concave, not linear \\
    P5 & larger influence in vacuum & no &
        compare vacuum and atmospheric events at equal $E$ \\
    P6 & one organising field & yes (in form) &
        thermal and mechanical channels decorrelate \\
    \bottomrule
  \end{tabular}
\end{table}

% =====================================================================
\section{Empirical calibration}
\label{sec:calibration}

\subsection{What is free, and what is a unit}

The framework has exactly one physically free parameter, the dimensionless
calibration constant $K$. The three reference scales $E_0$, $T_0$, $r_0$ are
choices of unit, not content: replacing them rescales $e$, $\tgrad$, $x$
simultaneously and leaves Eq.~\eqref{eq:theta_phys} invariant. We therefore
have one number to fix, and the rest of the theory is testable once it is.

\subsection{One anchor event suffices}

From any single event for which the source energy, the local thermal gradient
and an observed threshold radius are known, $K$ follows in closed form by
inverting Eq.~\eqref{eq:radius}:
\begin{equation}
  \widehat{K}
  \;=\; \frac{\tgrad\,\theta_c\,\ln\!\bigl(R_d/r_0\bigr)}{\sqrt{e}}
  \;=\; \frac{\theta_{\rm obs}\,\tgrad\,\ln x_{\rm obs}}{\sqrt{e}} .
  \label{eq:Kanchor}
\end{equation}
The second form applies when the field itself, rather than a material
threshold, is observed. Choosing one well-characterised high-energy event and
applying Eq.~\eqref{eq:Kanchor} removes the only free constant; it does not
confer any validity, and it is a one-parameter fit, so it cannot itself test
the framework.

\subsection{Many events: weighted least squares}

With $N$ events carrying measured field values
$\theta_i \pm \sigma_{\theta,i}$, define the model predictor
\begin{equation}
  p_i \;=\; \frac{\sqrt{e_i}}{\tgrad_i\,\ln x_i} .
  \label{eq:predictor}
\end{equation}
The model is then linear in $K$, $\theta_i = K\,p_i$, and the
weighted-least-squares estimator has the closed form
\begin{equation}
  \widehat{K} \;=\; \frac{\sum_i w_i p_i\theta_i}{\sum_i w_i p_i^{2}} ,
  \qquad
  \widehat{\sigma}_K = \left(\sum_i w_i p_i^{2}\right)^{-1/2},
  \qquad w_i = \sigma_{\theta,i}^{-2} .
  \label{eq:wls}
\end{equation}
Equations~\eqref{eq:predictor}--\eqref{eq:wls} are the estimator implemented
in the released package (with a non-negativity constraint $K\ge0$ and
covariance propagation from the nonlinear least-squares routine). Because the
model is linear in its only free parameter, no local minima exist and the fit
is unconditionally stable; the associated $\chi^{2}$,
\begin{equation}
  \chi^{2}(K) = \sum_i w_i\bigl(\theta_i - Kp_i\bigr)^{2},
  \qquad \chi^{2}_{\min} = \chi^{2}(\widehat{K}),
\end{equation}
with $N-1$ degrees of freedom, is the natural goodness-of-fit statistic once
$K$ has been absorbed. A correct model should give
$\chi^{2}_{\min}\approx N-1$; a systematically larger value is the
statistical signature of the wrong functional form, and is the quantitative
counterpart of the visual tests in P1 and P4.

\subsection{The two-step protocol}

We recommend a protocol that never uses $K$ to fit and $K$ to test at the same
time.
\begin{enumerate}
  \item \textbf{Anchor.} Choose one well-measured high-energy event, fix
        $E_0,T_0,r_0$ to convenient units, and solve Eq.~\eqref{eq:Kanchor}
        for $K$.
  \item \textbf{Parameter-free check.} Take a \emph{second} event at a
        different range and test Eq.~\eqref{eq:p2}. No fitted quantity enters
        this test. This is the step that can kill the theory.
  \item \textbf{Shape checks.} Over a decade of energies, test the linearity
        of $\ln R_d$ against $\sqrt{E}$ (P1) and the concavity of the
        field-versus-energy curve (P4).
  \item \textbf{Thermal check.} Sweep ambient temperature at fixed geometry
        and test $\thetaS\propto1/\tgrad$ (P3).
  \item \textbf{Publish the negative results.} A failed prediction is as
        informative as a successful one, and reporting it is what makes the
        framework a research programme rather than a claim.
\end{enumerate}

\subsection{What happens if P4 fails}

If the measured energy response is found to be $\propto E^{\alpha}$ with
$\alpha\neq 1/2$, the repair is transparent and worth stating in advance. Let
us generalise P-III to $\thetaS \propto e^{\alpha/2}$. Then, from
Eq.~\eqref{eq:alphagen},
\begin{equation}
  x_c = \exp\!\left(\frac{K\,e^{\alpha/2}}{\tgrad\,\theta_c}\right),
  \qquad
  E_c = E_0\left(\frac{\tgrad\,\theta_c\,\ln x}{K\,e^{\alpha/2}}\right)^{\!2},
  \label{eq:alphagen}
\end{equation}
and the parameter-free ratio of Eq.~\eqref{eq:p2} generalises to
\begin{equation}
  \frac{E_2}{E_1} = \left(\frac{\ln x_2}{\ln x_1}\right)^{\frac{2}{1+\alpha}} .
  \label{eq:alphap2}
\end{equation}
Two things follow. First, the parameter-freeness survives: Eq.~\eqref{eq:alphap2}
still contains neither $K$ nor $\theta_c$, so the discriminating experiment
of Sec.~\ref{sec:p2} remains available with $\alpha$ unknown. Second, the
exponent is measurable --- a multi-decade range-extension experiment measures
$\alpha$ directly. We regard the square-root response as a conjecture to be
measured, not a derived result, and Eq.~\eqref{eq:alphap2} is the
corresponding fallback. A related caveat applies to the threshold
$\theta_c$, for which no microphysical formula is offered: the framework
currently requires $\theta_c$ to be supplied as a material input.

\subsection{Status of the calibration to date}

No empirical calibration has been performed. In the released software $K$ is
set to unity and all reference scales to unity, which is a transparent
placeholder rather than a measurement; the figures in this paper are
illustrations of the model, not images of nature. Fig.~\ref{fig:field_decay}
shows the predicted slow fall-off, Fig.~\ref{fig:thermalmap} the
thermal--distance trade-off, Fig.~\ref{fig:radius} the exponential growth of
the destruction radius with $\sqrt{e}$ that embodies P1, and
Fig.~\ref{fig:breaking} the gentle quadratic-in-$\ln x$ growth of the
breaking energy that underlies P2. They are included so that a reader can see
the shapes the theory predicts; they carry no evidential weight.

\sugfigure{fig_field_decay.png}{field_decay}{Predicted field magnitude
$\lvert\thetaS\rvert$ against dimensionless distance $x=r/r_0$ at fixed
$\tgrad=1$, for several source energies $e$. The square-root energy response
(P-III) is visible in the unequal vertical spacing of the curves, and the
slow, long-reach logarithmic fall-off (P-II) in the flattening of every
curve. \emph{Illustration of the model, not a measurement.}}

\sugfigure{fig_thermal_energy_map.png}{thermalmap}{The thermal--distance
trade-off: contours of constant $\lvert\thetaS\rvert$ in the $(x,\tgrad)$ plane
at fixed $e$, i.e.\ the map implied by Eq.~\eqref{eq:thermalinv}. Moving closer
in $x$ can compensate for a larger thermal gradient, which is the
two-dimensional content of P-III and P-IV combined.
\emph{Illustration of the model, not a measurement.}}

\sugfigure{fig_vacuum_destruction_radius.png}{radius}{Predicted vacuum
destruction radius $R_d/r_0$ against dimensionless source energy $e$ on
logarithmic axes, for several thresholds $\theta_c$. Exponential growth of
reach with $\sqrt{e}$ is the embodiment of P1; note that a plot against
$\sqrt{e}$ on linear axes is a straight line, whereas the axes shown here are
equivalent up to the reparametrisation $e\to\sqrt{e}$.
\emph{Illustration of the model, not a measurement.}}

\sugfigure{fig_breaking_energy.png}{breaking}{Predicted breaking energy
$E_c/E_0$ against dimensionless distance $x$ on logarithmic axes, for several
thresholds. The gentle growth --- quadratic in $\ln x$ --- is the cheap-range
signature quantified by P2 in Eq.~\eqref{eq:p2}.
\emph{Illustration of the model, not a measurement.}}

% =====================================================================
\section{Conclusion}
\label{sec:conclusion}

SUG has been presented as a disciplined proposal: five explicit axioms, a
dimensionally honest central operator, a closed-form contour integral,
closed-form threshold inversions, and --- most importantly --- a set of
predictions each of which is paired with the route by which it could be
shown wrong. Its single most testable statement, the parameter-free energy
ratio of Eq.~\eqref{eq:p2}, requires no fitted parameter and can be applied
to a second event immediately.

Whether the framework survives contact with data is an open empirical
question, and the answer is not currently known. What this paper attempts to
establish is not that SUG is right, but that it is stated precisely enough
that the question can be settled decisively rather than argued about.

\subsection{Limitations and the validation agenda}
\label{sec:limitations}

We list the limitations explicitly, in the order in which they would have to
be addressed.

\begin{enumerate}
  \item \textbf{No calibration.} $K$ is free until fitted to a real anchor
        event. Only shapes and ratios are currently predictive; no absolute
        radius or energy is claimed.
  \item \textbf{Tension with the inverse-square law.} Postulate P-II, the
        central and most aggressive assumption, is in direct conceptual
        tension with the $1/r^{2}$ spread of radiative flux in three
        dimensions. It is not yet clear whether the SUG field is a rival
        description, a re-scaled effective quantity, or a bookkeeping
        device for an underlying derivation. This is the single most
        important open problem.
  \item \textbf{Source-core singularity.} The divergence at $x=1$ requires a
        physical prescription for $r_0$ and a cut-off or regularity
        condition. The interpretation as the edge of the source core is an
        assumption.
  \item \textbf{Square-root response asserted.} $\thetaS\propto\sqrt{E}$ is
        postulated, not derived. Section~\ref{sec:calibration} gives the
        $\alpha$-generalisation and shows that the key test survives it.
  \item \textbf{Material threshold.} $\theta_c$ has no microphysical
        formula and must be supplied per material or regime.
  \item \textbf{No data.} No experiment or validated simulation currently
        tests SUG. Every figure in this paper illustrates the model only.
\end{enumerate}

The honest summary is that SUG is a self-consistent axiomatic proposal that
has not yet confronted data. Its value at this stage is that it is cheap to
falsify: one calibration anchor plus one parameter-free cross-check would
either earn the framework a place or retire it.

% =====================================================================
\section*{Code Availability}
\label{sec:code}

All equations, operators, inversions, least-squares calibration and figures in
this paper are implemented in a released, tested, open-source Python package
\cite{sug2026}.

\begin{itemize}
  \item \textbf{Source repository (version control, issues, full derivation
        notes):}
        \url{https://github.com/sayan9168/Sayan-Universal-Gradient-Theory}
  \item \textbf{Archived, citable snapshot (Zenodo):}
        \url{https://doi.org/10.5281/zenodo.22723745}
        \quad (DOI \code{10.5281/zenodo.22723745})
  \item \textbf{PyPI package:} \code{pip install sayan-gradient}
        (MIT licence; Python $\ge 3.9$; requires \code{numpy} and
        \code{scipy}).
  \item \textbf{Interactive application:} a Streamlit web application
        (\code{streamlit run app.py}) exposes the $\Theta_S$ operator, the
        vacuum destruction radius, the energy breaking threshold, the $3$-D
        field surface and the P2 parameter-free ratio calculator.
\end{itemize}

The implementation is a direct transcription of the algebra in
Sec.~\ref{sec:framework}. The round-trip identities of
Sec.~\ref{sec:consistency}, the closed form of Eq.~\eqref{eq:contour} against
adaptive quadrature, and the parameter-free ratio of Eq.~\eqref{eq:p2} are all
asserted in an automated test suite, so that a reader can confirm that the
code and the printed equations agree. Reproducing the figures requires only
\code{python -m sug\_theory.plots --outdir figures} followed by
\code{pytest}.

% =====================================================================
\begin{thebibliography}{9}

\bibitem{buckingham1915}
E.~Buckingham, ``On the scaling of the equations of physics,''
\emph{Philosophical Magazine} \textbf{29}, 174 (1915).

\bibitem{clausius1865}
R.~Clausius, ``\"Uber die bewegende Kraft der W\"arme,''
\emph{Annalen der Physik} \textbf{126}, 353 (1865).

\bibitem{einstein1905}
A.~Einstein, ``Zur Elektrodynamik bewegter K\"orper,''
\emph{Annalen der Physik} \textbf{322}, 891 (1905).

\bibitem{sedov1959}
L.~I.~Sedov, \emph{Similarity and Dimensional Methods in Mechanics}
(Wiley, New York, 1959).

\bibitem{abramowitz1964}
M.~Abramowitz and I.~A.~Stegun (eds.),
\emph{Handbook of Mathematical Functions} (National Bureau of Standards,
Washington, 1964), Sec.~5.1 for the logarithmic integral $\Ei(\ln x)$.

\bibitem{sug2026}
S.~M., \emph{The Sayan Universal Gradient (SUG) Theory: A Proposed Framework
Coupling Energy Release, Thermal Gradients and Cosmic Distance}, GitHub
repository \code{sayan9168/Sayan-Universal-Gradient-Theory}, preprint version
2.1 (2026), \url{https://doi.org/10.5281/zenodo.22723745}.
SUG is an original proposal by the author; the works cited above provide
context and comparison only, and do not constitute endorsement or validation.

\end{thebibliography}

\end{document}
